\section{Objetivos generales}
\label{Objetivos generales}

\begin{itemize}
    
    
    \item \textbf{Desarrollar un sistema de entrenamiento virtual:} Este objetivo implica la creación de un sistema completo de entrenamiento virtual que integre herramientas de simulación avanzadas y algoritmos de aprendizaje automático. Se busca no solo diseñar escenarios virtuales realistas, sino también implementar un entorno interactivo que permita al robot militar aprender y adaptarse a diferentes situaciones de combate simuladas.
    
    
    \item \textbf{Construir un prototipo físico del robot militar:} Desarrollar un prototipo físico del robot basado en los resultados obtenidos de los entrenamientos virtuales. Este prototipo servirá como herramienta para validar y complementar los resultados virtuales con pruebas prácticas en condiciones controladas, permitiendo ajustes y mejoras en el diseño y rendimiento del robot.
    
    
    \item \textbf{Integrar simulación virtual con validación física para garantizar robustez:} Combina la simulación virtual con la validación física para asegurar la robustez y efectividad del sistema en situaciones reales de combate. Este enfoque integral garantiza que el robot esté completamente preparado para enfrentar los desafíos del campo de batalla, al haber sido entrenado y validado tanto en entornos virtuales como físicos.

    \item \textbf{ Describir un proceso de producción para un prototipo productivo:} Este objetivo implica detallar las etapas necesarias para la producción de un prototipo funcional. Incluye definir el diseño preliminar, seleccionar componentes, integrar subsistemas de hardware y software, desarrollar e implementar algoritmos para navegación y control, ensamblar el prototipo con componentes electrónicos y mecánicos.


    
\end{itemize}